\begin{textsample}{POS Dim 2 – vem\_ed\_22 – Score 4.00 – vem\_ed\_22/t114.txt}  \label{ex:f2_pos_178}
Fatecs em evento \textbf{on-line} de universidade indiana O evento \textbf{on-line} Student´s Research Conclave, organizado em 15 de dezembro de 2023 pelo Symbiosis Centre for Management Studies (Pune, Índia), selecionou trabalhos de 33 Instituições de Ensino Superior do Brasil (Fatecs / Centro Paula Souza), Índia, Catar e Zimbábue.
Marcaram \textbf{presença} quatro trabalhos das Fatecs Baixada Santista, Itaquaquecetuba e Pompeia – este último, derivou de estudo realizado sobre escassez de água no Nordeste brasileiro durante um Projeto Colaborativo Internacional (PCI / Cesu).
O trabalho da Fatec Baixada Santista \textbf{ficou} em segundo lugar geral e primeiro lugar na área de finanças, ganhando destaque no portal do Governo de São Paulo: https://www.saopaulo.sp.gov.br/ultimas-noticias/equipe-de-alunos-da-fatec-baixada-santista-e-vice-campea-em-evento-na-india/ Os \textbf{seguintes} artigos das Fatecs serão publicados no Annual Research Journal of SCMS 2024: Fatec Baixada Santista - Application of Principal Component Analysis (PCA) on a dataset of shares traded on the Brazilian Stock Exchange (B3) in the period 2018 – 2022: A feasibility study, dos estudantes Flávia Aparecida do Valle Escórcio; Gabriel Luiz dos Santos Silva e Thalita Carvalho Routh, orientados por Alexandre Garcia de Oliveira;

% matched lemmas: ficar, on-line, presença, seguinte
\end{textsample}
